 % --- Plik: chapter4_experiments.tex ---
 
 \chapter{Wyniki badań}
 \label{ch:wyniki} % To jest unikalny znacznik tego rozdziału
 
 Niniejszy rozdział służy prezentacji, eksperymentalnej walidacji oraz krytycznej analizie uzyskanych wyników badawczych, pomiarowych lub symulacyjnych. 
 
Treść powinna w sposób obiektywny dokumentować działanie i charakterystykę opracowanego rozwiązania, uwzględniając:
 \begin{itemize}
 	\item opis metodyki badawczej, warunków przeprowadzenia eksperymentu lub scenariuszy testowych,
 	
 	\item opracowane i czytelnie sformatowane dane wynikowe zaprezentowane w formie tabel zbiorczych oraz wykresów,
 	
 	\item interpretację uzyskanych zależności, analizę błędów pomiarowych oraz odniesienie wyników do założeń początkowych i stanu wiedzy.
 \end{itemize}
 \noindent \textbf{Stylistyka:} w tym rozdziale nie należy używać określeń subiektywnych (np. \textit{„wyniki są bardzo dobre”, „system działa świetnie”}). Należy posługiwać się wyłącznie językiem faktów, liczb i konkretnych wskaźników inżynierskich, stosując bezosobową formę czasu przeszłego.

